\begin{statementbox}
  \textbf{\textcolor{CStatement}{先想一个问题}}

  怎样才能让 $f(x)=ax^2+bx+c$ \textbf{无论 $x$ 取什么实数，结果都大于 $0$}？
  看图就明白：整条抛物线都要在 $x$ 轴\textbf{上方}，不能碰到它，更不能穿过它。

  \medskip
  \textbf{\textcolor{CStatement}{适用条件}}
  $a,b,c$ 为实数，$a\ne0$，并且要求对\textbf{所有实数} $x$ 成立。

  \medskip
  \textbf{\textcolor{CStatement}{核心结论（恒正）}}
  \[
    \boxed{\ f(x)>0\text{ 对一切 }x\in\mathbb R\text{ 成立}\ \Longleftrightarrow\ a>0\ \text{且}\ \Delta=b^2-4ac<0\ }
  \]
  \textbf{为什么是两个条件？} $a>0$ 管“开口朝上”；$\Delta<0$ 管“完全不碰 $x$ 轴”。两者缺一不可。

  \begin{center}
  \begin{tikzpicture}[x=0.95cm,y=0.79cm,>=stealth]
    \begin{scope}[shift={(0,0)}]
      \draw[->,gray] (-1.65,0)--(1.75,0) node[right] {$x$};
      \draw[->,gray] (0,-0.7)--(0,2.9) node[above] {$y$};
      \draw[very thick,green!55!black,domain=-1.38:1.38,samples=50] plot (\x,{\x*\x+0.55});
      \fill[green!55!black] (0,0.55) circle (1.8pt);
      \node[above] at (0,3.3) {\small\textbf{恒正}};
      \node[below,align=center] at (0,-0.83) {\scriptsize 开口向上，不相交\\[-1pt]\scriptsize $a>0,\ \Delta<0$};
    \end{scope}
    \begin{scope}[shift={(4.35,0)}]
      \draw[->,gray] (-1.65,0)--(1.75,0) node[right] {$x$};
      \draw[->,gray] (0,-0.7)--(0,2.9) node[above] {$y$};
      \draw[very thick,orange!85!black,domain=-1.58:1.58,samples=50] plot (\x,{\x*\x});
      \fill[orange!85!black] (0,0) circle (2pt);
      \node[above] at (0,3.3) {\small\textbf{仅非负}};
      \node[below,align=center] at (0,-0.83) {\scriptsize 只碰到一次 $x$ 轴\\[-1pt]\scriptsize $a>0,\ \Delta=0$};
    \end{scope}
    \begin{scope}[shift={(8.7,0)}]
      \draw[->,gray] (-1.65,0)--(1.75,0) node[right] {$x$};
      \draw[->,gray] (0,-0.7)--(0,2.9) node[above] {$y$};
      \draw[very thick,red!75!black,domain=-1.58:1.58,samples=50] plot (\x,{\x*\x-0.55});
      \fill[red!75!black] (-0.742,0) circle (1.6pt);
      \fill[red!75!black] (0.742,0) circle (1.6pt);
      \node[above] at (0,3.3) {\small\textbf{不恒正}};
      \node[below,align=center] at (0,-0.83) {\scriptsize 穿过 $x$ 轴\\[-1pt]\scriptsize $a>0,\ \Delta>0$};
    \end{scope}
  \end{tikzpicture}

  \end{center}

  \textbf{\textcolor{CStatement}{顺手记住（恒非负）}}
  \[
    \boxed{\ f(x)\ge0\text{ 对一切实数成立}\ \Longleftrightarrow\ a>0\ \text{且}\ \Delta\le0\ }
  \]
  \textbf{一字之差：}“大于 $0$”不许碰轴，“大于或等于 $0$”允许碰一下。
\end{statementbox}
